\section{Method}\label{sec:method}

The serial and the OpenMP builds are compiled from one C source file set. The serial build omits
\texttt{-fopenmp}, so every pragma is ignored and the same loops run on one thread. Any difference in time is
therefore due to parallel execution, not to a different algorithm.

\subsection{Pipeline}\label{sec:pipeline}
\textit{Stage 1, green channel.} $G(x,y)$ is taken from the RGB fundus image, where vessels have the highest
contrast. DRIVE images are LZW-compressed TIFF files, which we decode with our own C reader (verified
pixel-exact against Pillow); STARE images are PPM files.

\textit{Stage 2, CLAHE~\cite{zuiderveld1994}.} The image is split into an $8\times8$ grid of tiles of
$\lceil w/8\rceil\times\lceil h/8\rceil$ pixels. Each tile's 256-bin histogram is clipped at
$\kappa N_t/256$ counts ($N_t$ tile pixels, $\kappa$ the clip factor), the excess is spread uniformly over the
bins, and the cumulative histogram gives a lookup table. Each output pixel is a bilinear blend of the four
nearest tile tables.

\textit{Stage 3, multi-scale Frangi filter~\cite{frangi1998}.} With $I$ the CLAHE output scaled to $[0,1]$, the
second derivatives at scale $\sigma$ are computed by separable 1-D convolutions with sampled Gaussian kernels of
radius $r=\lceil3\sigma\rceil$: $I_{xx}=g_\sigma^{(y)}*(g''^{(x)}_\sigma*I)$,
$I_{yy}=g''^{(y)}_\sigma*(g^{(x)}_\sigma*I)$ and $I_{xy}=g'^{(y)}_\sigma*(g'^{(x)}_\sigma*I)$, with replicated
borders. The kernels are normalised so that $\sum g=1$, $\sum k\,g'(k)=-1$, $\sum g''=0$ and
$\sum \tfrac12k^2g''(k)=1$. One row sweep produces the three row-filtered images and one column sweep the three
derivatives, i.e.\ six 1-D passes per scale. With $a=\sigma^2I_{xx}$, $b=\sigma^2I_{xy}$ and
$d=\sigma^2I_{yy}$, the scale-normalised Hessian has eigenvalues
\begin{equation}
\lambda_{1,2}=\tfrac12(a+d)\mp\sqrt{\tfrac14(a-d)^2+b^2},
\end{equation}
relabelled so that $|\lambda_1|\le|\lambda_2|$. Vessels are dark on a bright background, so the response is
non-zero only when $\lambda_2>0$. With $R_B=|\lambda_1|/\lambda_2$ and $S^2=\lambda_1^2+\lambda_2^2$,
\begin{equation}
\mathcal V_\sigma=\exp\!\left(-\frac{R_B^2}{2\beta^2}\right)\left[1-\exp\!\left(-\frac{S^2}{2c^2}\right)\right],
\end{equation}
and $V(x,y)=\max_{\sigma\in\{1,2,3,4\}}\mathcal V_\sigma(x,y)$.

\textit{Stage 4, threshold and morphology.} $V$ is thresholded at a global value $\tau$ and multiplied by the
field-of-view (FOV) mask eroded by $e$ pixels, which removes the bright ring that the filter produces on the FOV
border. A $3\times3$ binary opening (erosion then dilation) removes speckle and a $3\times3$ closing bridges
one-pixel gaps.

\textit{Stage 5, evaluation.} True and false positives and negatives are counted inside the (uneroded) FOV and
converted to Dice $=2TP/(2TP+FP+FN)$, sensitivity, specificity and accuracy.

\subsection{Parameter selection}\label{sec:tuning}
All parameters are chosen on the 20 DRIVE training images only, by maximising the mean per-image Dice over a
grid: $\kappa\in\{1,2,3,4\}$, $\beta\in\{0.5,1\}$,
$c\in\{0.005,\allowbreak 0.01,\allowbreak 0.02,\allowbreak 0.04,\allowbreak 0.08,\allowbreak 0.16\}$,
$e\in\{4,8,12\}$ and 40 thresholds spaced logarithmically in $[10^{-3},0.9]$ (5{,}760 settings). The selected
setting is then fixed and applied unchanged to the 20 DRIVE test images, scored against the first observer's
manual segmentation, and to the 20 STARE images, scored against Hoover's labels~\cite{hoover2000}. STARE has no
FOV masks, so the FOV is the set of pixels whose maximum channel value exceeds 35, cleaned by a $3\times3$
opening.

\subsection{OpenMP design}\label{sec:omp}
Table~\ref{tab:design} lists the parallel decomposition of each stage. All loops use
\texttt{schedule(static)} unless stated, and each output element is computed by exactly one thread, so
floating-point summation order is the same as in the serial build. Three alternative designs are kept as
ablations: a dynamic schedule for the CLAHE tile loop, a pixel-parallel CLAHE histogram with
\texttt{\#pragma omp atomic} increments, and a Frangi decomposition that runs the four scales concurrently with
one thread each instead of parallelising inside each scale.

\begin{table}[t]
\centering
\caption{Parallel decomposition by stage.}\label{tab:design}
\footnotesize
\begin{tabular}{@{}p{0.23\columnwidth}p{0.70\columnwidth}@{}}
\toprule
Stage & OpenMP decomposition \\
\midrule
Green channel & \texttt{parallel for} over pixels \\
CLAHE & \texttt{parallel for} over the 64 tiles, each thread owning its tile histogram and lookup table; \texttt{parallel for} over rows for the bilinear blend \\
Gaussian derivatives & row sweep: \texttt{parallel for} over rows with a thread-private padded row buffer; column sweep: \texttt{parallel for} over output rows, inner loop over $x$ vectorised \\
Eigen + vesselness & \texttt{parallel for} over pixels, fused with the running maximum over scales \\
Threshold & \texttt{parallel for} over pixels, fused with the FOV mask \\
Morphology & four $3\times3$ passes, each a \texttt{parallel for} over rows \\
Metrics & \texttt{reduction(+:TP,FP,TN,FN)} \\
\bottomrule
\end{tabular}
\end{table}

\subsection{Correctness checks}\label{sec:verify}
For every test image the serial build writes the CLAHE output, the vesselness map $V$ and the final mask, and
every OpenMP configuration (each design variant at 1, 2, 4 and 8 threads) writes the same files. We report the
fraction of identical CLAHE bytes, identical mask pixels and bit-identical $V$ values, and the maximum
$|\Delta V|$. The comparison is repeated for builds with floating-point contraction disabled
(\texttt{-ffp-contract=off}). As an external check, $V$ from the serial build is compared with
\texttt{skimage.\allowbreak filters.\allowbreak frangi}~\cite{vanderwalt2014} applied to the same CLAHE output with the same $\sigma$,
$\beta$ and $c$; the two implementations normalise derivatives differently, so we report correlation rather
than equality.

\subsection{Platform and measurement protocol}\label{sec:protocol}
All timed runs use one Kaggle CPU notebook: an Intel Xeon at 2.20\,GHz exposing 4 logical CPUs (2 physical
cores with 2-way Hyper-Threading~\cite{marr2002}), 32\,KiB L1d and 256\,KiB L2 per core, a 55\,MiB L3
and 31\,GiB RAM, with a cgroup quota of 4 CPUs, running Linux 6.18 with GCC 13.3 (\texttt{-O3 -march=native}, OpenMP 4.5). Nothing is
timed on the authors' laptops. Thread counts are 1, 2, 4 and 8 (the last oversubscribes the 4 logical CPUs);
threads are bound with \texttt{OMP\_PROC\_BIND=close} and \texttt{OMP\_PLACES=threads} unless stated. Each
configuration runs 5 untimed warm-up iterations and 30 timed iterations per image over all 20 DRIVE test images.
Time is measured with \texttt{omp\_get\_wtime} (OpenMP build) or the monotonic clock (serial build) around each
stage; file I/O and image decoding are excluded. Speedup is $S(p)=T_\text{serial}/T(p)$ with
$T_\text{serial}$ from the pragma-free build, efficiency is $E=S/p$ and the Karp--Flatt serial
fraction~\cite{karp1990} is $e=(1/S-1/p)/(1-1/p)$. Throughput is images per second for one image at a time.

For resolution scaling, three test images are bilinearly resampled to $2\times$, $4\times$ and $8\times$ linear
size (0.33 to 21.1 megapixels) before timing, with the FOV mask resampled by nearest neighbour and the erosion
radius scaled accordingly; the number of timed iterations is reduced at large sizes (10, 10, 5 and 3). The
three design ablations are timed on the first five test images with the same protocol. A thread-placement
experiment repeats the sweep on the same five images with one thread per physical core first
(\texttt{OMP\_PLACES=cores}, \texttt{OMP\_PROC\_BIND=spread}) and with binding disabled. STARE is timed on all 20
images with 2 warm-up and 10 timed iterations.
